\scanpage{98}
Soit $K$ une extension \deleted{\unsure{gal}} \unsure{tot|tot.} réelle finie de $Q$,
$L$ ext. quadratique tot. imaginaire de $K$,
$q=p^{\unsure{n|-}}$ une puissance de nb premier. On cherche
les $\lambda\in K$ \uline{entiers} tels que $\lambda\bar\lambda=q^i$, i.e.
$|\lambda|=q^{i/2}$, \inserted{pour $i$ convenable $\in\mathbb N$.} Ce sont des unités exc. $p$, et
ils engendrent \deleted{le groupe des} un sous-groupe du
groupe \deleted{\unsure{$\forall$} des unités exc. $p$, \unsure{Soit les $\lambda$ entiers}}
formé des $\lambda$ unités exc. $p$ tels que
$\lambda\bar\lambda=q^i$ ($i\in\mathbb Z$). Notons que \deleted{ce sont des
unités \uline{absolues} exc. $p$. Donc finalement il en est
ces $\lambda$ s'écrivent $q^{i/2}u$, avec $u$} si $q$ est un
carré de \deleted{\unsure{$Q$|$K$}}$K$, on peut écrire un $\lambda$
sous la forme $\lambda=q^{i/2}u$, où $u$ est une unité
exc. $p$ telle que $u\bar u=1$, \unsure{Ce|En} qui implique que
$u$ est une \uline{unité absolue} exc. $p$. Les $u$ forment
un groupe $U$, et l'hom
\[
U\unsure{\longrightarrow|-}\mathbb Z^P
\qquad u\unsure{\mapsto|\sim}\bigl(v_{\mathfrak p}\unsure{(u)|N(u)}\bigr)_{\mathfrak p\in P}
\]
(où $P$ désigne l'ens des $\mathfrak p$ de $L$ sur $p$) est
de noyau fini. D'après la prop. plus haut,
son image est \inserted{\unsure{contenue} et} d'indice fini \unsure{du|dans} $\mathbb Z^{\deleted{\unsure{Q}}P'}$, où
\deleted{$P'$ est l'ens des $\mathfrak p$} formé des diviseurs de trace
\inserted{(rel. à $K$)} nulle. Donc si pour $\forall$ $\mathfrak p_0$ \deleted{\unsure{$\in$}\illegible{1}} idéal premier de
$K$ sur $p$ qui \uline{\unsure{splitte|splitt} en deux} dans $L$, on choisit
un $\mathfrak p$ au dessus, on trouve un isom à un
groupe fini près $U\simeq\mathbb Z^Q$, $Q$ étant \unsure{l'ens|l'un} des $\mathfrak p_0$
% archive sheet number: 97

